% arara: pdflatex: { draft: yes } unless exists(toFile('tesis.aux'))
% arara: biber if changed (toFile('bibliografia.bib')) || found('log', 'Citation')
% arara: --> || found ('log', 'Please \\(re\\)run Biber')
% arara: pdflatex: { synctex: yes } until !found('log', '\\(?(R|r)e\\)?run (to get|LaTeX)')

\documentclass{tesisFC}

\usepackage{mycats}
\usepackage[backend=biber,defernumbers=true,style=alphabetic,giveninits=true]{biblatex}
\usepackage[style=mexican]{csquotes}
\usepackage{mathtools} %Extensión y corrección de errores de amsmath
\usepackage{amssymb}
\usepackage{enumitem}
\usepackage[predefined={parent=chapter}]{keytheorems} %Para ambientes de teoremas
\usepackage{tikz}
\usetikzlibrary{shapes,arrows,spy,positioning,snakes}
%\usepackage[colorlinks, allcolors=blue]{hyperref}
\usepackage{pdfcomment}
\hypersetup{colorlinks, allcolors=blue}

%Para tener axiomas como teoremas numerados
\AtBeginDocument{%
  \newkeytheorem{axioma}[sibling=theorem]
}

\addbibresource{bibliografia.bib}
\nocite{*}

%Funciones especiales
\DeclareMathOperator{\sub}{Sub}
\DeclareMathOperator{\homs}{Hom}
\DeclareMathOperator{\seno}{sen}

\title{La categoría de Conjuntos como subcategoría de la categoría de Espacios}
\author{Andrés Esteban Escudero Robles}
\grado{Matemático}
\date{2026}
\tipo{tesis}
\tutorW{tutor}
\tutor{Grado y nombre}

%\includeonly{2espacios}

\begin{document}
\frontmatter
\portadatesis
\tableofcontents*


\mainmatter
\include{0intro}
\chapter{Geometría Diferencial Sintética}
\label{geodif}

En esta sección usaremos los axiomas que expone Mclarty \cite{mclarty1988defining} para trabajar
rigurosamente con infinitesimales y fundamentar al cálculo, así como al estudio de variedades.
\begin{axioma}[label=ax:gds1]
Existe un espacio, $R$, que cuenta con una estructura de anillo conmutativo con unidad.
\end{axioma}
\begin{axioma}[label=ax:gds2]
El anillo $R$ no es trivial, es decir, $1\not = 0$.
\end{axioma}
\begin{axioma}[label=ax:gds3]
En la lógica interna de $R$, $R$ es un campo.
\end{axioma}
Y el axioma de Kock-Lawvere \cite{lawvere1979categorical}, para el cual es necesario definir
al espacio de infinitesimales $D$, el cual será el igualador de $0\colon R\to R$, el mapa constante $0$, y $(-)^2\colon R\to R$.
Equivalentemente, $D$ es el subespacio (subobjeto) de $R$, $D= \{d \in R\colon d^2=0 \}$.
\begin{axioma}[label=ax:gds4]
$(\forall f \in R^D)(\exists ! b \in R)(\forall d \in D) (f(d) = f(0)+d \cdot b)$.
\end{axioma}
Intuitivamente, \ref{ax:gds4} dice que toda función de $D$ a $R$ es lineal. A continuación 
mostramos algunas propiedades de los infinitesimales y que estos cuatro axiomas ya nos 
permiten hacer cálculo básico.

\section{El espacio de infinitesimales}
Los infinitesimales con los que trabajamos en GDS son números $d\in D$ que al cuadrado dan cero, y por supuesto,
esto implica que todas las potencias mayores a dos también dan cero, pues $d^n=d^2\cdot d^{n-2}=0$ para $n>2$. Esto
corresponde con los argumentos de Newton, Leibniz, Euler, entre otros, donde se eliminan los términos de orden
superior de una cantidad infinitesimal. Sin embargo, si queremos que $R$ sea un campo y tenemos que $\lnot(d=0)$,
$d$ es invertible, por lo que $d^2=0$ implica $d=0$, entonces sucede que $\lnot\lnot(d=0)$, pues estamos usando
lógica constructivista. Por otro lado, si para todo $d\in D$, $d=0$ y fijamos $f\colon  D \to R$, obtenemos que
$f(0)=f(d)=f(0)+d \cdot b$ aplicando \ref{ax:gds4}, pero como $d=0$, $f(d)=f(0)+d \cdot 0=f(0)+d \cdot 1$, y por unicidad
de $b$, $b=0=1$, lo cual es inconsistente con \ref{ax:gds2}. Así que también $\lnot (\forall d \in D)(d=0)$ es verdad.
De esta manera, $R$ puede ser un campo a pesar de los infinitesimales, en el sentido de que $\lnot(r=0)\implies \exists s\in R(s\cdot r=1)$.

\section{Infinitesimales y Microcancelación}
Si resulta que $(\forall d \in D) (d\cdot r = d\cdot s)$ con $r,s\in R$, podemos concluir
$r=s$. Demostración: Definimos $f\colon  D\to R$ como $f(d)=d\cdot r = d\cdot s$, luego,
por \ref{ax:gds4}, $f(d) = f(0)+d\cdot b$ donde $f(0) = 0\cdot r = 0$, entonces $d\cdot b=d\cdot r = d\cdot s$
con $b$ única. Por lo tanto, $r=b=s$.

\section{Operaciones con infinitesimales}
Si $r\in R$ y $d\in D$, entonces $d\cdot r \in D$, porque $(d\cdot r)^2=d^2\cdot r^2=0\cdot r^2=0$.
Si $d_1\in D$ y $d_2\in D$, entonces $(d_1+d_2)^2=d_1^2+2\cdot d_1\cdot d_2+d_2^2=2\cdot d_1\cdot d_2$.

\section{La derivada y reglas básicas}
Sea $f\colon  R\to R$, consideramos a $g\colon  D\to R$ tal que $g(d)=f(a+d)$ para alguna
$a\in R$ fija. Por \ref{ax:gds4}, $f(a+d)=g(d)=g(0)+d\cdot b$, pero $g(0)=f(a+0)=f(a)$, entonces
$f(a+d)=f(a)+d\cdot b$ y definimos a $b$ como la derivada de $f$ en $a$, $f'(a)$. Dado que $a$ era
arbitrario, en realidad definimos a una función $f'\colon  R \to R$, la derivada de $f$, que
puede volver a ser diferenciada con el mismo procedimiento. Concluimos que todas las funciones
de $R$ en $R$ son suaves.

\begin{definition}[label=def:derivada]
  Sea $f\colon  R\to R$ y $a\in R$. Definimos $f'(a)$ como el único número tal que para todo $d\in D$
  \begin{equation*}
    f(a+d)=f(a)+d\cdot f'(a)
  \end{equation*}
  Notemos que si pudiéramos dividir entre $d$, tendríamos $f'(a) = \frac{f(a+d)-f(a)}{d}$, lo cual es la definición clásica de derivada.
\end{definition}

% {\color{red}Borrar esto después. Ahora se puede hacer referencia a lo que
% definas o enuncies, ver definición~\ref{def:defrivada} o el teorema~\ref{thm:reglasderivada}.}

\begin{theorem}[label=thm:reglasderivada]
  Si $f\colon  R\to R$,\ $g\colon  R\to R$,\ $a,c\in R$ y $d\in D$, entonces
\begin{enumerate}[label=\textbf{\alph*)}]
    \item $c' = 0$
    \item $(c\cdot f)' =c\cdot f'$
    \item $\id' = 1$
    \item $(f+g)'=f'+g'$
    \item $(f\circ g)' =(f'\circ g)\cdot g'$
    \item $(f\cdot g)' =f'\cdot g+f\cdot g'$
\end{enumerate}
\end{theorem}
\begin{proof}
  Probemos d) como ejemplo. $(f+g)(a+d) = f(a)+g(a)+d\cdot (f+g)'(a)$
  por definición de derivada, pero también $(f+g)(a+d) = f(a+d)+g(a+d)=f(a)+d\cdot f'(a)+g(a)
  +d\cdot g'(a)$. Cancelando términos comunes obtenemos $d\cdot (f+g)'(a) = d \cdot[f'(a)+g'(a)]$
  y por microcancelación, $(f+g)'(a)=f'(a)+g'(a)$ con $a$ cualquiera, entonces
  $(f+g)'=f'+g'$.
  
  Para demostrar a) consideremos
  \begin{align*}
   c(a+d) &= c(a)+d\cdot c'(a)\\
    c &= c+d\cdot c'(a)\\
    d\cdot 0 &= d\cdot c'(a)\\
    c'(a) &= 0\\
    c' &= 0
  \end{align*}

  Ahora, para mostrar el inciso b), tenemos que
  \begin{align*}
     (c\cdot f)(a+d) &= c\cdot f(a)+d\cdot (c\cdot f)'(a)\\
      f(a+d) &= f(a)+d\cdot f'(a)\\
      c\cdot f(a+d) &= c\cdot f(a)+ c\cdot d\cdot f'(a)\\
      d\cdot [c\cdot f'(a)] &= d\cdot (c\cdot f)'(a)\\
      (c\cdot f)'(a) &=c\cdot f'(a)\\
      (c\cdot f)' &=c\cdot f' 
  \end{align*}

  El inciso c) se sigue de:
  \begin{align*}
     \id (a+d) &= \id (a)+d\cdot \id '(a)\\
      a+d &= a+d\cdot \id '(a)\\
      d\cdot 1 &= d\cdot \id '(a)\\
      \id '(a) &= 1\\
      \id ' &= 1\\
  \end{align*}

  Para el inciso e) hacemos lo siguiente
  \begin{align*}
     (f\circ g)(a+d) &= f(g(a))+d\cdot (f\circ g)'(a)\\
      (f\circ g)(a+d) &= f(g(a)+d\cdot g'(a))\\
      (f\circ g)(a+d) &= f(g(a))+[d\cdot g'(a)] \cdot f'(g(a))\\
      d\cdot (f\circ g)'(a) &= d\cdot [f'(g(a))\cdot g'(a)]\\
      (f\circ g)'&=(f'\circ g)\cdot g'\\
  \end{align*}

  Finalmente, para el inciso f) tenemos
  \begin{align*}
     (f\cdot g)(a+d) &= f(a)\cdot g(a)+d\cdot (f\cdot g)'(a)\\
      (f\cdot g)(a+d) &= f(a+d)\cdot g(a+d)\\
      (f\cdot g)(a+d) &= [f(a)+d\cdot f'(a)]\cdot [g(a)+d\cdot g'(a)]\\
      (f\cdot g)(a+d) &= f(a)\cdot g(a)+d\cdot[f'(a)\cdot g(a)+g'(a)\cdot f(a)]+d^2\cdot f'(a)\cdot g'(a)\\
      d\cdot (f\cdot g)'(a)&=d\cdot[f'(a)\cdot g(a)+f(a)\cdot g'(a)]\\
      (f\cdot g)'&=f\cdot g+f\cdot g' \qedhere
  \end{align*}
\end{proof}

\section{Ejemplos}
Sea $f(x) = 5x^2+3$. En ese caso, $f(x+d)=5(x^2+2xd+d^2)+3=5x^2+3+df'(x)$, entonces $d\cdot 10x=d\cdot f'(x)$.
Por tanto, $f'(x)=10x$.

Con axiomas más fuertes podemos introducir números transcendentales y por ejemplo, a la función seno. Para encontrar
la derivada de $seno(x)$, debemos evaluar $seno(a+d)=seno(a)\cdot cos(d)+cos(a)\cdot seno(d)$ con $a\in R$ y $d\in D$. Entonces queremos
saber cuanto valen $cos(d)$ y $seno(d)$, así que hay que considerar un triángulo rectángulo con un ángulo tan pequeño
que vale $d$. Pero a escalas infinitesimales, el círculo se ve como una línea recta, es decir, la longitud de arco para
trazar el ángulo $d$ es tan pequeña que el arco en realidad es una recta y por tanto, si trazamos el triángulo en un círculo unitario, el cateto opuesto al ángulo $d$ también mide $d$.
Ahora, si la hipotenusa es de longitud 1, lo anterior implica que $seno(d)=d$ y $cos(d)=\sqrt{1^2-d^2}=1$. Por lo tanto, $seno(a+d)=seno(a)+d\cdot cos(a)$ y $seno'(a)=cos(a)$
para toda $a\in R$ por \ref{def:derivada}.

\begin{center}
  \begin{tikzpicture}
    \draw (1,1) -- (5,1) -- (5,3) -- (1,1);
    \draw (5,1.5) -- (4.5,1.5) -- (4.5,1);
    \draw (1.8,1) -- (1.8,1.4);

    \node at (2,1.2) {$d$};
    \node at (3,2.2) {$1$};
    \node at (5.2,2) {$d$};
  \end{tikzpicture}
\end{center}

\section{Más axiomas para Espacios}

\begin{definition}[label=def:discreto]
  Un espacio $E$ es discreto si y solo si es decidible, es decir, $(x=y) \lor \lnot(x = y)$ es verdad con $x,y$ variables
  variando sobre $E$.
\end{definition}

Mclarty le pide más axiomas a Espacios.

\begin{axioma}[label=ax:gds5]
Existe un espacio $N$ que es un objeto de números naturales.
\end{axioma}

Se puede demostrar que $N$ es es un espacio discreto.

\begin{definition}[label=puntos]
Los puntos de un espacio $E$ son sus elementos globales, morfismos de la forma $1\to E$
\end{definition}

\begin{axioma}[label=ax:gds6]
Para todo espacio $E$, $E \cong  0 $ o $E$ tiene al menos un punto.
\end{axioma}

Esto implica que si $b_1, b_2\in B$ entonces $(b_1=b_2)\lor \lnot(b_1=b_2)$ es verdad incluso si $B$ no es discreto.

\begin{axioma}[label=ax:gds7]
Todo espacio $E$ tiene un subespacio $\Gamma E$ discreto tal que todo punto de $E$ está en $\Gamma E$. 
\end{axioma} 
\chapter{El topos Espacios}
\label{espacios}

(Equivalencia de ser discreto)

\begin{proposition}
  Los subespacios de un espacio discreto son discretos
\end{proposition}
\begin{proof}
  Sea $A$ un espacio discreto y $m\colon S\to A$ un subespacio de $A$, queremos demostrar que la 
  diagonal $\Delta_S\colon S\to S\times S$ es un subobjeto complementado. Por hipótesis, para $\Delta_A\colon A\to A\times A$
  existe su complemento, $C$, tal que $C\cap A=0$ y $C\cup A=A\times A$. Notamos que
  $m$ induce $m\times m\colon S\times S \rightarrowtail A\times A$ y $(m\times m)^{-1}(\Delta_A)=\Delta_S$ ya que el siguiente es un producto fibrado:
  \begin{equation*}
    \begin{tikzcd}
        %T\ar[bend left=30]{rrd}\ar[bend right=30]{rdd}[swap]{t}
          %\ar[dashed]{rd}{\pi_S t}\\[-2ex]
        S\ar{r}{m} \ar[tail]{d}[swap]{\Delta_S} & A\ar[tail]{d}{\Delta_A}\\
        S\times S\ar{r}[swap]{m\times m} & A\times A.
    \end{tikzcd}
  \end{equation*}
  Para demostrar que el diagrama es un producto fibrado necesitamos que el cuadrado conmute
  y que para cualquier objeto $T$ con flechas que haga conmutar al diagrama exterior, exista 
  un único morfismo $T\to S$ que haga conmutar los triángulos. Como $\Delta_A m = \langle id_A, id_A\rangle m = m\times m$ la conmutatividad
  del cuadrado es inmediata. Para un $T$ con flechas $t\colon T\to S\times S$, $u\colon T\to A$ tales
  que $\Delta_A u = \langle m, m\rangle t$, proponemos a $\pi_S t$ como la flecha de $T$
  a $S$, aprovechando que $t=\langle\pi_S t, \pi_S t\rangle$ donde $\pi_S\colon S\times S \to S$
  es la proyección a $S$. De esta manera tenemos la conmutatividad del triángulo izquierdo
  y, junto con la conmutatividad del cuadrado exterior e interior, obtenemos que $\Delta_A u
  = \Delta_A m \pi_S t$, pero $\Delta_A$ es mono, entonces $u=m \pi_S t$ y el triángulo superior
  también conmuta. La unicidad de $\pi_S t$ es inmediata ya que $\Delta_S$ es
  mono.\pdfcomment{Tal vez se pueda quitar la demostración}
  

  Dado que $S\times S$ es un subobjeto de $A\times A$ a través de $m\times m$, la imagen
  inversa $(m\times m)^{-1}$ es un morfismo de álgebras de Heyting $Sub(A\times A)\to Sub(S\times S)$,
  con lo cual $(m\times m)^{-1}$ preserva complementos y por lo anterior $\Delta_S\colon S\to S\times S$ es un
  subobjeto complementado con complemento $(m\times m)^{-1}(C)$.

\end{proof}

\begin{proposition}[label=prop:terminaldiscreto]
  El objeto terminal de un topos es un objeto decidible.
\end{proposition}
\begin{proof}
  Queremos mostrar que $\Delta_1\colon 1\to 1\times 1$ es un subobjeto complementado, pero $1\times 1 \cong 1$
  entonces $1$ es el subobjeto maximal, por lo que $0 \to 1\times 1$ es el complemento que buscamos.
\end{proof}

\begin{theorem}
  En un topos cualquiera, si $A\in \sub(1)$ y $A$ tiene un punto $a\colon 1\to A$, entonces $A \cong 1$.
\end{theorem}
\begin{proof}
  Sea $A$ un subespacio de $1$, con $m\colon A\to 1$ monomorfismo. Juntando nuestras hipótesis obtenemos al siguiente diagrama:
  \begin{equation*}
    \begin{tikzcd}
      1\ar{r}{a}\ar[swap]{rrd}{\text{id}} & A\ar{r}{m} & 1\ar[shift left = 0.3ex, no head]{d}{}\ar[shift right = 0.3ex, no head]{d}{} \\
       & {} & 1.
    \end{tikzcd}
  \end{equation*}
  El diagrama conmuta pues acaba en el terminal y como la composición es igual a la identidad
  en $1$, $m$ es un epimorfismo que se escinde aparte de ser mono. Por tanto, $m$ es un isomorfismo y $A\cong 1$.
\end{proof}

\begin{corollary}[label=cor:bivaluado]
  El clasificador de subobjetos de Espacios tiene dos puntos.
\end{corollary}
\begin{proof}
  Sea $\Omega$ el clasificador de subobjetos de Espacios, entonces $\sub(1)\cong \homs(1,\Omega)$
  y por el teorema anterior, $\sub(1)=\{0,1\}$, así que $\homs(1,\Omega)=\{0,1\}$.
\end{proof}

\begin{definition}
  Un objeto de números naturales (ONN) en una categoría con objeto terminal, es una tercia $(N,o,s)$
  donde $N$ es un objeto de la categoría, $o\colon 1\to N$ y $s\colon N\to N$ son morfismos tales que para todo
  objeto $A$ y morfismos $a\colon 1\to A$ y $g\colon A\to A$ existe un único morfismo $f\colon N\to A$ tal que
  los siguientes diagramas conmutan:
  \begin{equation*}
    \begin{tikzcd}
      1\ar{r}{o}\ar[swap]{rd}{a} & N\ar{d}{f}\ar{r}{s} & N\ar{d}{f} \\
      & A\ar[swap]{r}{g} & A.
    \end{tikzcd}
  \end{equation*}
  Intuitivamente, la flecha $o$ señala al cero y $s$ es la función sucesor, por lo que
  $f$ es la función definida recursivamente tal que $f(0)=a$ y $f(n+1)=g(f(n))$.
\end{definition}

En una categoría cartesiana cerrada, un ONN permite también hacer recursión con parámetros,
es decir, definir funciones usando que $f(a,0)=g(a)$ y $f(a,n+1)=h(a,f(a,n))$ \cite{johnstone2002sketches}.

\begin{proposition}[label=prop:onncoproducto]
  Sea $(N,o,s)$ un ONN en una categoría $\mathcal{C}$ con coproductos binarios, entonces
  \begin{equation*}
    \begin{tikzcd}
      1\ar{r}{o} & N & N,\ar[swap]{l}{s}
    \end{tikzcd}
  \end{equation*}
  es un coproducto.
\end{proposition}
\begin{proof}
  La demostración se resume en que el ONN es el objeto inicial de la categoría de T-álgebras con
  $T\colon \mathcal{C}\to \mathcal{C}$ el endofuntor $1+(-)$. Dada una T-álgebra $(A,a)$, con $a\colon 1+A\to A$,
  por definición de coproducto tenemos $i_1\colon 1\to 1+A$ e $i_A\colon A\to
  1+A$, entonces existe una única $f\colon N\to A$ que hace conmutar al siguiente diagrama:
  \begin{equation*}
    \begin{tikzcd}
      1\ar{r}{o}\ar[swap]{rd}{ai_1} & N\ar[dashed]{d}{f}\ar{r}{s} & N\ar[dashed]{d}{f} \\
      & A\ar[swap]{r}{ai_A} & A.
    \end{tikzcd}
  \end{equation*}
  Para que $f$ realmente sea un morfismo entre $(N,[o,s]\colon 1+N\to N)$ y $(A,a\colon 1+A\to A)$ necesitamos que $f[o,s]=aT(f)$.
  \begin{equation*}
    \begin{tikzcd}
      1\ar{r}{i_1}\ar[swap]{rdd}{ai_1}\ar{rd}{o} & 1+N\ar{d}{[o,s]} & N\ar[swap]{l}{i_N}\ar{ldd}{fs}\ar{ld}{s} & {} & 1\ar{r}{i_1}\ar[swap]{rdd}{ai_1}\ar{rd}{i_1} & 1+N\ar[swap]{d}{Tf} & N\ar[swap]{l}{i_N}\ar{ldd}{fs}\ar[swap]{ld}{i_Af} \\
        & N\ar{d}{f} & {} & {} & {} & 1+A\ar{d}{a} & {} \\
        & A & {} & {} & {} & A. & {}
    \end{tikzcd}
  \end{equation*}
  La construcción de $f$ nos permite concluir que los diagramas son conmutativos, así que
  por la propiedad universal del coproducto obtenemos lo que buscábamos y por tanto, $(N,[o,s])$ es el
  objeto inicial de las T-álgebras. Pero el morfismo estructural de una T-álgebra inicial es un isomorfismo,
  gracias a que $(TN, T[o,s])$ es también una T-álgebra y no es difícil verificar que el morfismo inducido $g\colon N\to TN$ es
  inverso por ambos lados de $[o,s]$ (en general, del morfismo estructural de una T-álgebra inicial \cite{johnstone2002sketches}).
  De esta manera, $N\cong 1+N$.
\end{proof}

\begin{definition}[label=def:predecesor]
  En una categoría con ONN y productos binarios, es posible definir al morfismo predecesor
  como el único $p\colon N\to N$ que hace conmutar al siguiente diagrama:
  \begin{equation*}
    \begin{tikzcd}
      1\ar{r}{o}\ar[swap]{rd}{o\times o} & N\ar[dashed]{d}{(id,p)}\ar{r}{s} & N\ar[dashed]{d}{(id,p)} \\
      & N\times N\ar[swap]{r}{(s\pi_1,\pi_1)} & N\times N.
    \end{tikzcd}
  \end{equation*}
  Intuitivamente, la conmutatividad del diagrama nos dice que $p(0)=0$ y $p(n+1)=n$.
\end{definition}

\begin{proposition}[label=prop:sucesormono]
  Sea $(N,o,s)$ un ONN en una categoría con productos binarios, entonces $s$ es un monomorfismo.
\end{proposition}
\begin{proof}
  Sea $p\colon N\to N$ el predecesor de~\ref{def:predecesor}, notemos en la construcción de $p$ que
  la conmutatividad en la segunda entrada dice que $ps=\pi_1(id_N,p)=id_N$. Por tanto,
  $s$ tiene retracto y es mono.
\end{proof}

\begin{proposition}[label=prop:complementosucesor]
  Sea $(N,o,s)$ un ONN en un topos, entonces $o$ es el complemento de $s$ en $\sub(N)$.
\end{proposition}
\begin{proof}
  Ya vimos que $s$ es mono en \ref{prop:sucesormono} y $o$ es trivialmente mono. Por \ref{prop:onncoproducto}
  $N\cong 1+N$ pero en un topos los coproductos son disjuntos, y por tanto $o$ es el complemento de $s$.
\end{proof}

\begin{proposition}
  En toda categoría con ONN podemos hacer ``inducción'', es decir, si $m\colon N'\to N$ es un subobjeto
  tal que tanto $o$ como la restricción de $s$ a $N'$ se factorizan a través de $N'$, entonces $N'\cong N$.
\end{proposition}
\begin{proof}
  Sean $o'\colon 1\to N'$ y $s'\colon N'\to N'$ las factorizaciones, por la propiedad universal del ONN, existe un único 
  morfismo $f\colon N\to N'$ tal que el siguiente diagrama es conmutativo.
  \begin{equation*}
    \begin{tikzcd}
      1\ar{r}{o}\ar{rd}{o'}\ar[swap]{rdd}{o} & N\ar{d}{f}\ar{r}{s} & N\ar{d}{f} \\
      & N'\ar[swap]{r}{s'}\ar{d}{m} & N'\ar{d}{m} \\
      & N\ar{r}{s} & N.
    \end{tikzcd}
  \end{equation*}
  Por lo que $mf\colon N\to N'\to N$ tiene que ser la identidad en $N$, pero esto implica que $m$ además de ser mono tiene
  inverso por la derecha, así que $m$ es un isomorfismo.
\end{proof}

\begin{theorem}
  Sea $\mathcal{E}$ un topos con y $(N,o,s)$ un objeto de números naturales, entonces $N$ es un objeto decidible.
\end{theorem}
\begin{proof}
  El plan es ver a $\Delta\colon N\to N\times N$ como intersección de subobjetos complementados en $\sub(N\times N)$ (así va la demostración de Johnstone pero hay una más directa con la lógica interna).
  Primero, en un topos tenemos recursividad con parámetros que nos permite definir a la resta truncada como 
  $\dot{-}\colon N\times N\to N$ tal que $\dot{-}(a,0)=a$ y $\dot{-}(a,n+1)=p(\dot{-}(a,n))$ donde $p$ es el predecesor de~\ref{def:predecesor}.
  Ahora construimos a los subobjetos:
  \begin{equation*}
    \begin{tikzcd}
      N_1\ar{r}{}\ar{d}{} & 1\ar{d}{o} \\
      N\times N\ar{r}{\dot{-}} & N.
    \end{tikzcd}
  \end{equation*}
  Que son complementados gracias a que ser complementado es estable bajo productos fibrados. Como $n\dot{-}n=0$ para toda $n$,
  $\Delta$ es subobjeto de $N_1$ y por lo tanto $\Delta$ también tiene complemento en $\sub(N\times N)$.
\end{proof}

\begin{definition}
  Sea $M$ la categoría de espacios discretos y los morfismos entre ellos en Espacios. $M$ es una subcategoría plena de Espacios.
\end{definition}

\begin{theorem}[label=thm:funtordiscretos]
  \ref{ax:gds6} y \ref{ax:gds7} inducen una pareja de funtores adjuntos entre $M$ y Espacios.
\end{theorem}
\begin{proof}
  Para cada $A$ un espacio tenemos un monomorfismo $c_A\colon  \Gamma A \to A$ dado
  que $\Gamma A$ es un subobjeto. Veamos que $c_A$ es couniversal hacia
  $A$ desde $\Delta\colon  M \to Espacios$, la inclusión de $M$. Sea $B$ un objeto de $M$ y $f\colon  B \to A$,
  deseamos encontrar un morfismo que haga conmutar al siguiente diagrama:
  \begin{equation*}
    \begin{tikzcd}
      \Delta B \ar{r} \ar[swap]{rd}{f} & \Delta\Gamma A \ar[tail]{d}{c_A} \\
                                & A.
    \end{tikzcd}
  \end{equation*}
  Notemos que si $B$ fuera el objeto inicial ya acabamos. En otro caso, 
  consideramos al producto fibrado entre $f$ y $c_A$. Además, por \ref{ax:gds6}, podemos
  tomar $b\colon  1\to B$.
  \begin{equation*}
    \begin{tikzcd}
        1\ar[bend left=30]{rrd}\ar[bend right=30]{rdd}[swap]{b}
          \ar[dashed]{rd}{b}\\[-2ex]
        &[-2ex] f^{-1} (\Gamma A)\ar{r}{\overline{f}} \ar[tail]{d}[swap]{c_B} & \Gamma A\ar[tail]{d}{c_A}\\
        & B\ar{r}[swap]{f} & A.
    \end{tikzcd}
  \end{equation*}
  Para ver que el exterior conmuta recordemos que según \ref{ax:gds7}, los puntos de $A$ se factorizan a
  través de $\Gamma A$ y $fb$ es un punto de $A$. Entonces los puntos de $B$ se factorizan
  a través de $f^{-1} (\Gamma A)$. Pero $f^{-1} (\Gamma A)$ es discreto ya que es subobjeto
  de $B$, que es discreto, por lo que $\Gamma B = f^{-1} (\Gamma A)$ y dado que $B$ está
  en $M$, $\Gamma B = B$. Así que $\overline{f}$ es el morfismo que cumple lo que queríamos
  y además es único gracias a que $c_A$ es mono. Por tanto, $\Gamma : Espacios \to M$ es
  un funtor que manda a los espacios a sus espacios discretos asociados
  y $\Gamma$ es adjunto derecho de $\Delta$ con $c_\square : \Delta\Gamma \to \text{Id}$ la counidad.
\end{proof}

\begin{theorem}
  Sea $\mathcal{E}$ un topos y $C=(G,\epsilon,\delta)$ una comonada en $\mathcal{E}$ tal que
  $G$ es exacto por la izquierda. Entonces la categoría de coálgebras $\mathcal{E}_C$ es un topos.
\end{theorem}
\begin{lemma}
  La categoría de coálgebras $\mathcal{E}_C$ tiene límites finitos.
\end{lemma}
\begin{proof}
  Es suficiente mostrar que la categoría de coálgebras tiene productos e igualadores finitos. 
  Como $\mathcal{E}$ tiene productos y $G$ es funtor, dadas $(A,s),(B,t)$ coálgebras de $C$ tenemos:
  \begin{equation*}
    \begin{tikzcd}
        GA & GA\times GB \ar[swap]{l}{\pi_{GA}} \ar{r}{\pi_{GB}} & GB\\
        & G(A\times B). \ar{ul}{G(\pi_A)} \ar[dashed]{u}{\phi} \ar[swap]{ur}{G(\pi_B)}
    \end{tikzcd}
  \end{equation*}
  Con $\phi$ natural en los objetos A y B.En este caso, gracias a que $G$ es izquierdo exacto, $\phi$ es un isomorfismo. Proponemos
  $(A\times B, \phi^{-1}(s\times t))$ como el producto de $(A,s)$ y $(B,t)$. Veamos que
  es una coálgebra de $C$.
  \begin{equation*}
    \begin{tikzcd}
        A\times B\ar{r}{s\times t}\ar[swap]{d}{\pi_A}\ar[phantom]{2-2}{(1)} & GA\times GB \ar[swap]{d}{\pi_{GA}}\ar{r}{\phi^{-1}}\ar[phantom]{2-3}[near start]{(2)} 
          & G(A\times B)\ar{dl}{G(\pi_A)}\ar{d}{\epsilon_{A\times B}} \\
        A\ar{r}{s}\ar[swap]{dr}{id_A} & GA\ar{d}{\epsilon_A}\ar[phantom]{r}{(4)}\ar[phantom]{3-1}[near start]{(3)} & A\times B \ar{dl}{\pi_A} \\
        {} & A.
    \end{tikzcd}
  \end{equation*}
  En el diagrama (1) conmuta por la definición de $s\times t$, (2) conmuta gracias a la construcción
  de $\phi$, (3) conmuta ya que $(A,s)$ es coálgebra y (4) conmuta pues es un cuadrado de naturalidad
  de $\epsilon$. Por tanto $\pi_A=\pi_A \epsilon_{A\times B} \phi^{-1}(s\times t)$ y dado que podemos
  repetir el argumento para $\pi_B$, concluimos que $\epsilon_{A\times B} \phi^{-1}(s\times t)$
  es la identidad en $A\times B$. Esto es la ecuación de la counidad, entonces seguimos con la
  coasociatividad. De manera similar a $\phi$ podemos encontrar un isomorfismo $\psi\colon G(GA\times GB) \to G^2A\times G^2B$,
  lo cual nos da el siguiente diagrama de isomorfismos.
  \begin{equation*}
    \begin{tikzcd}
      G(GA\times GB)\ar{r}{G(\phi^{-1})} & G^2(A\times B) \\
      G^2A\times G^2B.\ar{u}{\psi^{-1}}\ar[swap]{ur}{G(\phi^{-1})\psi^{-1}}
    \end{tikzcd}
  \end{equation*}
  Para la coasociatividad queremos ver la conmutatividad del siguiente diagrama:
  \begin{equation*}
    \begin{tikzcd}
      A\times B\ar[phantom]{rd}{(1)}\ar{r}{s\times t}\ar[swap]{d}{s\times t} & GA\times GB\ar{r}{\phi^{-1}}\ar{d}{\delta_A\times \delta_B}
        & G(A\times B)\ar{dd}{\delta_{A\times B}} \\
      GA\times GB\ar{r}{G(s)\times G(t)}\ar[swap]{d}{\phi^{-1}} & G^2A\times G^2B\ar{d}{\psi^{-1}} \\
      G(A\times B)\ar[phantom]{ur}{(2)}\ar[swap]{r}{G(s\times t)} & G(GA\times GB)\ar[phantom]{uur}{(3)}\ar{r}{G(\phi^{-1})} & G^2(A\times B).
    \end{tikzcd}
  \end{equation*}
  En (1) sabemos que el diagrama en cada coordenada conmuta pues $(A,s),(B,t)$ son coálgebras. 
  Para (2) tenemos:
  \begin{equation*}
    \begin{tikzcd}
      GA\ar[swap]{d}{G(s)} & GA\times GB\ar[swap]{l}{\pi_{GA}}\ar{r}{\pi_{GB}}\ar[dashed]{d}{G(s)\times G(t)} & GB\ar{d}{G(t)}\\
      G^2A & G^2A\times G^2B\ar[swap]{l}{\pi_{G^2A}}\ar{r}{\pi_{G^2B}} & G^2B \\
      G^2A\ar{u}{id_{G^2A}} & G(GA\times GB)\ar[swap]{l}{G(\pi_{GA})}\ar{u}{\psi}\ar{r}{G(\pi_{GB})} & G^2B\ar[swap]{u}{id_{G^2B}} \\
      GA\ar{u}{G(s)} & G(A\times B)\ar[swap]{l}{G(\pi_A)}\ar{u}{G(s\times t)}\ar{r}{G(\pi_B)} & GB\ar[swap]{u}{G(t)} \\
      & GA\times GB.\ar{ul}{\pi_{GA}}\ar{u}{\phi^{-1}}\ar[swap]{ur}{\pi_{GB}}
    \end{tikzcd}
  \end{equation*}
  Cada parte conmuta gracias a cómo construimos los morfismos usando la propiedad universal del producto,
  por lo que $G(s)\times G(t)=\langle G(s)\pi_{GA},G(t)\pi_{GB} \rangle = \psi G(s\times t)\phi^{-1}$,
  lo cual nos da la identidad deseada al poscomponer con $\psi^{-1}$. En el caso de (3) observamos que:
  \begin{equation*}
    \begin{tikzcd}
      {} & GB\ar{r}{\delta_B} & G^2B & {} \\
      GA\times GB\ar{ur}{\pi_{GB}}\ar{r}{\phi^{-1}}\ar[swap]{dr}{\pi_A} & G(A\times B)\ar[swap]{u}{G(\pi_B)}\ar{r}{\delta_{A\times B}}\ar{d}{G(\pi_A)}
        & G^2(A\times B)\ar{u}{G^2(\pi_B)}\ar{r}{\psi G(\phi)}\ar[swap]{d}{G^2(\pi_A)}
        & G^2A\times G^2GB\ar[swap]{ul}{\pi_{G^2B}}\ar{dl}{\pi_{G^2A}} \\
      & GA\ar[swap]{r}{\delta_A} & G^2A.
    \end{tikzcd}
  \end{equation*}
  Donde todo conmuta ya que son diagramas de productos y cuadrados de naturalidad de $\delta$, así
  que $\delta_A\times\delta_B=\psi G(\phi)\delta_{A\times B}\phi^{-1}$, pero $\psi G(\phi)$ es un
  isomorfismo, entonces obtenemos la igualdad que queríamos.

  Ahora consideramos a los morfismos de coálgebras $f\colon (X,u)\to (A,s)$ y $g\colon (X,u)\to (B,t)$ para demostrar
  la propiedad universal de $(A\times B, \phi^{-1}(s\times t))$.
  \begin{equation*}
    \begin{tikzcd}
      X\ar[swap]{d}{u}\ar{r}{f\times g} & A\times B\ar{dr}{\phi^{-1}(s\times t)}\ar{d}{s\times t} \\
      GX\ar[swap]{r}{G(f)\times G(g)} & GA\times GB\ar[swap]{r}{\phi^{-1}} & G(A\times B).
    \end{tikzcd}
  \end{equation*}
  El cuadro conmuta pues $f$ y $g$ son morfismos de coálgebras, además como $G$ preserva
  productos tenemos que $\phi^{-1}(G(f)\times G(g))=G(f\times g)$, lo cual demuestra la
  existencia del morfismo $f\times g\colon (X,u)\to (A\times B,\phi^{-1}(s\times t))$. La unicidad
  se da gracias a la unicidad de $f\times g$ en $\mathcal{E}$. Por tanto $\mathcal{E}_C$
  tiene productos finitos y procedemos a mostrar que tiene igualadores. Consideramos $(A,s)$,
  $(B,t)$ coálgebras y $f,g\colon (A,s)\to (B,t)$. Proponemos $(E,u)$ como sigue:
  \begin{equation*}
    \begin{tikzcd}
      E\ar{r}{e}\ar[swap,dashed]{d}{u} & A\ar[shift left=0.75ex]{r}{f}\ar[swap, shift right=0.75ex]{r}{g}\ar{d}{s} & B\ar{d}{t} \\
      GE\ar[swap]{r}{G(e)} & GA\ar[shift left=0.75ex]{r}{G(f)}\ar[swap, shift right=0.75ex]{r}{G(g)} & GB. \\
    \end{tikzcd}
  \end{equation*}
  Los cuadros conmutan ya que $f$ y $g$ son morfismos de coálgebras, entonces $s\circ e$ iguala
  a $G(f)$ con $G(g)$, lo cual nos da la existencia de $u$. $GE$ es el igualador de $G(f)$ y $G(g)$ salvo isomorfismo
  porque $G$ preserva límites. Veamos que $(E,u)$ es coálgebra, para ello empezamos con la ``counidad'',
  es decir, queremos demostrar la conmutatividad del siguiente diagrama:
  \begin{equation*}
    \begin{tikzcd}
      E\ar{r}{u}\ar[swap]{dr}{id_E} & GE\ar{d}{\epsilon_E} \\
      & E
    \end{tikzcd}
  \end{equation*}
  Notemos que $\epsilon_AG(e)$ iguala a $f,g$ por la naturalidad de $\epsilon$.
  \begin{equation*}
    \begin{tikzcd}
      E\ar{r}{e} & A\ar[shift left=0.75ex]{r}{f}\ar[swap, shift right=0.75ex]{r}{g} & B & {} E\ar{r}{e}\ar[swap]{d}{u} & A\ar{r}{id_A}\ar{d}{s} & A \\
      GE\ar[dashed]{u}{\epsilon_E}\ar[swap]{r}{G(e)} & GA\ar[shift left=0.75ex]{r}{G(f)}\ar[swap, shift right=0.75ex]{r}{G(g)}\ar{u}{\epsilon_A} & GB\ar{u}{\epsilon_B} & {} GE\ar[swap]{r}{G(e)} & GA\ar[swap]{ur}{\epsilon_A} \\
      E.\ar{u}{u}
    \end{tikzcd}
  \end{equation*}
  La flecha inducida de $GE$ a $E$ debe ser $\epsilon_E$, de nuevo por la naturalidad de $\epsilon$.
  El diagrama de la derecha conmuta por la definición de $u$ y porque $s$ es la estructura de coálgebra de $A$,
  entonces $\epsilon_AG(e)u=e$, así que, por la propiedad universal de $E$ como igualador, $\epsilon_Eu=id_E$. Para la coasociatividad, queremos ver que el siguiente diagrama conmuta:
  \begin{equation*}
    \begin{tikzcd}
      E\ar{r}{u}\ar[swap]{d}{u} & GE\ar{d}{\delta_E} \\
      GE\ar[swap]{r}{G(u)} & G^2E.
    \end{tikzcd}
  \end{equation*}
  Para ello consideramos el siguiente diagrama:
  \begin{equation*}
    \begin{tikzcd}
      A\ar{r}{s} & GA\ar[shift left=0.75ex]{r}{G(s)}\ar[swap, shift right=0.75ex]{r}{\delta_A} & G^2A \\
      E\ar{u}{e}\ar[swap]{r}{u} & GE\ar[shift left=0.75ex]{r}{G(u)}\ar[swap, shift right=0.75ex]{r}{\delta_E}\ar{u}{G(e)} & G^2E.\ar[swap]{u}{G^2(e)}
    \end{tikzcd}
  \end{equation*}
  El cuadrado de la izquierda conmuta por la construcción de $u$ y ambos cuadrados derechos conmutan, en un caso
  porque $G$ es funtor y en el otro caso porque $\delta$ es una transformación natural. Más aún, $\delta_As=G(s)s$
  por la coasociatividad de $A$, pero $G^2(e)$ es un igualador y como los cuadros conmutan, $\delta_Ase$, $G(s)se$
  también igualan a $G^2(f)$ con $G^2(g)$, entonces por la propiedad universal de $G^2E$ obtenemos que $\delta_Eu=G(u)u$. Ahora veamos que $(E,u)$ es el igualador de $f$ y $g$ en $\mathcal{E}_C$.
  \begin{equation*}
    \begin{tikzcd}
      E\ar{r}{e}\ar[swap]{d}{u} & A\ar[shift left=0.75ex]{r}{f}\ar[swap, shift right=0.75ex]{r}{g}\ar{d}{s} & B\ar{d}{t} \\
      GE\ar[swap]{r}{G(e)} & GA\ar[shift left=0.75ex]{r}{G(f)}\ar[swap, shift right=0.75ex]{r}{G(g)} & GB.
    \end{tikzcd}
  \end{equation*}
  Claramente $e$, visto como morfismo de coálgebras, iguala a $f$ con $g$. Sea $h\colon (X,w)\to (A,s)$ otro morfismo
  de coálgebras que iguala a $f$ y a $g$, entonces por la propiedad universal de $E$ y $GE$ tenemos:
  \begin{equation*}
    \begin{tikzcd}
      X\ar[dashed]{r}{j}\ar[swap]{d}{w} & E\ar{d}{u} \\
      GX\ar[swap,dashed]{r}{k} & GE.
    \end{tikzcd}
  \end{equation*}
  Por construcción, $ej=h$ entonces $G(e)G(j)=G(h)$, lo cual implica que $G(j)=k$. Como $G(h)$ iguala a $G(f)$ con $G(g)$,
  $G(h)w\colon X\to GA$ induce una única flecha de $X$ a $GE$.
  \begin{equation*}
    \begin{tikzcd}
      GE\ar{r}{G(e)} & GA & GE\ar{rr}{G(e)} & {} & GA\\
      GX\ar{u}{G(j)}\ar{ur}{G(h)} & {} & E\ar{u}{u}\ar{r}{e} & A\ar{ur}{s} \\
      X\ar{u}{w}\ar[swap]{uur}{G(h)w} & {} & X\ar{u}{j}\ar[swap]{ur}{h}\ar[swap]{rr}{w} & {} & GX.\ar[swap]{uu}{G(h)}
    \end{tikzcd}
  \end{equation*}
  El diagrama de la izquierda es la definición de $k=G(j)$. El diagrama de la derecha conmuta por la definición
  de $u$, la definición de $j$ y porque $h$ es un morfismo de coálgebras. Por lo tanto, $uj=G(j)w$ y existe
  $j\colon (X,w)\to (E,u)$. La unicidad de $j$ similarmente se da gracias a la unicidad de $j$ en $\mathcal{E}$. De esta manera, $\mathcal{E}_C$ tiene límites finitos.
\end{proof}

\begin{lemma}
  La categoría $\mathcal{E}_C$ tiene exponenciales.
\end{lemma}
\begin{proof}
  Seguiremos la demostración de Mac Lane \cite{maclane2012sheaves}. Sean $(A,s), (B,t)$ y $(C,u)$ coálgebras, entonces tenemos la siguiente cadena de isomorfismos naturales:
  \begin{equation*}
    \homs_{\mathcal{E}}(A\times B, C) \cong \homs_{\mathcal{E}}(U(A,s), C^B) \cong \homs_{\mathcal{E}_C}((A,s), R(C^B)).
  \end{equation*}
  Recordando que tenemos una adjunción entre el funtor olvidadizo, $U$, y el funtor colibre, $R$.
  La coálgebra colibre de un objeto $X$ es $(GX, \delta_X)$ y un morfismo $h\colon U(Y,\alpha)\to X$ es enviado a
  $G(h)\alpha\colon (Y,\alpha)\to RX$ bajo la adjunción. Por el isomorfismo anterior, un buen candidato
  para el exponencial $(C,u)^{(B,t)}$ es $R(C^B)$, sin embargo, debemos restringirnos a los morfismos
  de coálgebras, por lo que sospechamos que hay un subobjeto $(E,v)$ de $R(C^B)$ tal que tengamos
  el siguiente isomorfismo natural:
  \begin{equation*}
    \homs_{\mathcal{E}_C}((A,s)\times (B,t), (C,u)) \cong \homs_{\mathcal{E}_C}((A,s), (E,v)).
  \end{equation*}
  Lo anterior implicaría que $(E,v)$ es el exponencial que buscamos. Sea $f\colon (A,s)\times (B,t) \to (C,u)$,
  definimos a $f'$ como su transpuesta exponencial en $\mathcal{E}$ y sea $f''=G(f')s$. Si traducimos
  la condición de que $f$ sea morfismo de coálgebras en una condición sobre $f''$, sabremos como
  definir a $(E,v)$. Recordando que para cualesquiera $X,Y$ en $\mathcal{E}$ hay un isomorfismo 
  $\phi\colon G(X\times Y)\to GX\times GY$ natural en $X,Y$ y utilizando la construcción del producto en el lema anterior, la condición de que $f$ sea un morfismo de 
  coálgebras se expresa con el siguiente diagrama:
  \begin{equation*}
    \begin{tikzcd}
      A\times B\ar{r}{s\times t}\ar[swap]{d}{f} & GA\times GB\ar{r}{\phi^{-1}} & G(A\times B)\ar{d}{G(f)} \\
      C\ar{rr}{u} & {} & GC.
    \end{tikzcd}
  \end{equation*}
  Ahora transponemos cada morfismo empezando por $G(f)\phi^{-1}(s\times t)$ con la siguiente observación:
  \begin{equation*}
    \begin{tikzcd}
      A\times B\ar[swap]{d}{=}\ar{r}{s\times 1_B}\ar[phantom]{rd}{(1)} & GA\times B\ar[swap]{d}{Gf'\times 1_B}\ar{r}{1_{GA}\times t}\ar[phantom]{rd}{(2)} & GA\times GB\ar[swap]{d}{Gf'\times 1_{GB}}\ar{r}{\cong}\ar[phantom]{rd}{(3)} & G(A\times B)\ar[swap]{d}{G(f'\times 1_B)}\ar{r}{Gf}\ar[phantom]{rd}{(4)} & GC\ar{d}{=}\\
      A\times B\ar[swap]{r}{f''\times 1_B}              & G(C^B)\times B\ar[swap]{r}{1\times t}                       & G(C^B)\times GB\ar[swap]{r}{\phi^{-1}}                 & G(C^B\times B)\ar[swap]{r}{G(ev_{B,C})}             & GC.
    \end{tikzcd}
  \end{equation*}
  Donde la fila de arriba es precisamente $G(f)\phi^{-1}(s\times t)$ y cada cuadro conmuta, en (1) por la definición de $f''$, (2) es inmediato, (3) es por la 
  naturalidad de $\phi$ y (4) conmuta pues es aplicar al funtor $G$ a la definición de $f'$ como 
  transpuesta. La conmutatividad de (4), luego (3), etcétera, nos permite concluir que la fila de
  arriba es igual a la fila de abajo, que es un morfismo más sencillo de transponer.
  \begin{equation*}
    \overline{G(ev)\phi^{-1}(1\times t)(f''\times 1_B)} = \overline{G(ev)\phi^{-1}(1\times t)}f'' = (GC)^t\overline{G(ev)\phi^{-1}}f''.
  \end{equation*}
  Por otro lado, $\overline{uf} = u^Bf'$. Si nombramos $p=\overline{G(ev)\phi^{-1}}$, que $f$ sea
  morfismo de coálgebras es equivalente a que $u^Bf' = (GC)^tpf''$. Ahora pasamos la igualdad a
  través de la adjunción.
  \begin{align*}
    G(u^Bf')s &= G(u^B)G(f')s = G(u^B)f'' \\
    G((GC)^tpf'')s &= G((GC)^t)G(p)G(f'')s = G((GC)^t)G(p)\delta_{C^B}f''.
  \end{align*}
  Recordando que $f''\colon (A,s)\to (G(C^B),\delta_{C^B})$. Como $f''$ iguala a los anteriores
  morfismos entre coálgebras colibres, aprovechamos que $\mathcal{E}_C$ tiene límites finitos
  y proponemos a $(E,v)$ como el igualador de $G(u^B)$ con $G((GC)^t)G(p)\delta_{C^B}$.
  \begin{equation*}
    \begin{tikzcd}
      (E,v)\ar{r}{i} & (G(C^B),\delta_{C^B})\ar{r}{G(u^B)}\ar{d}{\delta_{C^B}} & (G((GC)^B),\delta_{(GC)^B}) \\
      (A,s)\ar[dashed]{u}{f^*} & (G(G(C^B)),\delta_{G(C^B)})\ar{r}{G(p)} & (G((GC)^{GB}),\delta_{(GC)^{GB}}).\ar{u}{G((GC)^t)} \\
    \end{tikzcd}
  \end{equation*}
  El mapeo de $f$ a $f^*$ define un isomorfismo natural ya que cada asignación se hizo a través
  de una propiedad universal, por lo que $(E,v)\cong (C,u)^{(B,t)}$.
\end{proof}
Procedemos a mostrar que $\mathcal{E}_C$ tiene un clasificador de subobjetos, para lo cual necesitaremos dos lemas.
\begin{lemma}
  Sea (A,s) una coálgebra y $m\colon D\rightarrowtail A$ un subobjeto de $A$ en $\mathcal{E}$. Entonces hay
  a lo más un morfismo $d\colon D\to GD$ tal que $(D,d)$ sea coálgebra y $m$ un morfismo de coálgebras.
\end{lemma}
\begin{proof}
  Supongamos que existe un $d\colon D\to GD$ tal que $(D,d)$ sea coálgebra y $m$ un morfismo de coálgebras. Entonces tenemos
  el siguiente diagrama:
  \begin{equation*}
    \label{eq:subalgebra}
    \begin{tikzcd}
      D\ar[tail]{r}{m}\ar[swap]{d}{d} & A\ar{d}{s} \\
      GD\ar[swap]{r}{G(m)} & GA.
    \end{tikzcd}
  \end{equation*}
  Como $m$ es mono y $G$ es exacto por la izquierda, $G(m)$ es mono, y por lo tanto, $d$ debe ser único.
\end{proof}
\begin{lemma}
  Si $(A,s)$ es coálgebra, $m\colon D\to A$ es un monomorfismo y $d\colon D\to GD$ cumple que $G(m)d=sm$, $(D,d)$ 
  es coálgebra y el diagrama~\ref{eq:subalgebra} es un producto fibrado en $\mathcal{E}$.
\end{lemma}
\begin{proof}
  Veamos que $(D,d)$ es coálgebra.
  \begin{equation*}
    \begin{tikzcd}
      D\ar{rrr}{d}\ar[swap]{ddd}{d}\ar{rd}{m} & {} & {} & GD\ar{ddd}{Gd}\ar{ld}{Gm} \\
      {} & A\ar{r}{s}\ar{d}{s} & GA\ar{d}{Gs} & {} \\
      {} & GA\ar{r}{\delta_A} & G^2A & {} \\
      GD\ar{rrr}{\delta_D}\ar{ru}{Gm} & {} & {} & G^2D.\ar{lu}{G^2m}
    \end{tikzcd}
  \end{equation*}
  El cuadrado interior conmuta dado que $(A,s)$ es coálgebra, el trapecio inferior conmuta por la
  naturalidad de $\delta$ y los demás trapecios conmutan por la hipótesis de que $G(m)d=sm$. Para que
  el exterior conmute, basta ver que $G^2(m)G(d)d=G^2(m)\delta_Dd$, ya que $G^2(m)$ es mono, entonces
  tenemos lo siguiente:
  \begin{align*}
    G^2(m)G(d)d &= G(s)G(m)d = G(s)sm \\
    G^2(m)\delta_Dd &= \delta_A G(m)d = \delta_Asm.
  \end{align*}
  Y como el interior conmuta, concluimos que el exterior también y por tanto, $d$ cumple la coasociatividad.
  Para la counidad, consideremos el siguiente diagrama:
  \begin{equation*}
    \begin{tikzcd}
      D\ar{r}{d}\ar[swap,tail]{d}{m} & GD\ar{r}{\epsilon_D}\ar[tail]{d}{Gm} & D\ar[tail]{d}{m} \\
      A\ar{r}{s} & GA\ar{r}{\epsilon_A} & A.
    \end{tikzcd}
  \end{equation*}
  El cuadrado de la izquierda conmuta por hipótesis y el de la derecha por la naturalidad de $\epsilon$,
  esto basta para mostrar que el exterior conmuta. Como $(A,s)$ es coálgebra, $\epsilon_As=id_A$, entonces 
  $m\epsilon_Dd=id_Am=m id_D$, pero $m$ es mono así que $\epsilon_Dd=id_D$ y $d$ cumple la counidad.
  Observemos con más cuidado el diagrama anterior. Supongamos que tenemos un objeto $X$ y morfismos $p\colon X\to GD$ y $q\colon X\to A$ tales que $G(m)p=sq$.
  \begin{equation*}
    \begin{tikzcd}
      X\ar[bend left=30]{rrd}{p}\ar[bend right=30]{rdd}[swap]{q}
          \ar[dashed]{rd}{\epsilon_Dp}\\[-2ex]
      & [-2ex] D\ar{r}{d}\ar[swap]{d}{m} & GD\ar{d}{Gm} \\
      & A\ar[swap]{r}{s} & GA.
    \end{tikzcd}
  \end{equation*}
  Notemos que:
  \begin{align*}
    G(m)d\epsilon_Dp = sm\epsilon_Dp = s\epsilon_AG(m)p &= s\epsilon_Asq = sq = G(m)p \\
    m\epsilon_Dp = \epsilon_AG(m)p &= \epsilon_Asq = q.
  \end{align*}
  Dado que $G(m)$ es mono, nuestra propuesta $\epsilon_Dp$ hace conmutar a los triángulos deseados
  y es único porque $m$ es mono. Por tanto, el diagrama es un producto fibrado.
\end{proof}
\begin{lemma}
  La categoría $\mathcal{E}_C$ tiene un clasificador de subobjetos.
\end{lemma}
\begin{proof}
  Empezamos con $v\colon 1\to \Omega$ el clasificador de subobjetos en $\mathcal{E}$ y definimos a $\tau\colon G\Omega\to \Omega$
  como el morfismo que clasifica al mono $G(v)\colon G1\to G\Omega$. Sean $(A,s),(D,d)$ coálgebras con $m\colon (D,d)\rightarrowtail (A,s)$
  un monomorfismo y sea $h\colon A\to \Omega$ el morfismo que clasifica a $m$ en $\mathcal{E}$. Entonces
  tenemos el siguiente diagrama:
  \begin{equation*}
    \label{eq:clasificadorm}
    \begin{tikzcd}
      D\ar{r}{d}\ar[swap]{d}{m} & GD\ar{r}{G(!_D)}\ar{d}{Gm} & G1\ar{r}{!_{G1}}\ar{d}{Gv} & 1\ar{d}{v} \\
      A\ar{r}{s} & GA\ar{r}{Gh} & G\Omega\ar{r}{\tau} & \Omega. \\
    \end{tikzcd}
  \end{equation*}
  Donde no solo cada cuadro conmuta, sino que es un producto fibrado. El cuadrado de la izquierda debido
  al lema anterior, el cuadrado de en medio es $G$ aplicado a la construcción de $h$, pero $G$ preserva
  productos fibrados, y el cuadrado de la derecha es producto fibrado por la definición de $\tau$.
  Por lo tanto, el exterior es un producto fibrado y $\tau G(h)s$ clasifica a $m$. Entonces $h=\tau G(h)s$,
  Donde $G(h)s$ es justamente la transposición de $h$ bajo la adjunción $U\dashv R$. Más aún
  \begin{align*}
    id_{G\Omega}G(h)s = G(\tau G(h)s) s &= G(\tau)G^2(h)G(s)s \\
    &=G(\tau)G^2(h)\delta_A s  \tag{Coasociatividad de s} \\
    &=G(\tau)\delta_{\Omega}G(h)s \tag{Naturalidad de $\delta$}.
  \end{align*}
  Es decir, $G(h)s$ iguala a $id_{G\Omega}$ con $G(\tau)\delta_{\Omega}$, que son morfismos de coálgebras,
  así que $G(h)s$ se factoriza a través del igualador
  \begin{equation*}
    \begin{tikzcd}
      (\Omega_G,u)\ar{r}{e} & (G\Omega,\delta_{\Omega})\ar[shift left=0.75ex]{r}{id_{G\Omega}}\ar[swap, shift right=0.75ex]{r}{G(\tau)\delta_{\Omega}} & (G\Omega,\delta_{\Omega}) \\
      (A,s).\ar[dashed]{u}{h^*}\ar[swap]{ru}{G(h)s}
    \end{tikzcd}
  \end{equation*}
  Esta factorización se da para cualquier monomorfismo de coálgebras, en particular, como $G$
  es exacto por la izquierda, hay una coálgebra trivial $(1,\psi)$ con el isomorfismo $\psi\colon 1\to G1$ y
  una sub-coálgebra $!\colon (1,\psi)\rightarrowtail (1,\psi)$ cuyo morfismo característico en $\mathcal{E}$ es $v\colon 1\to \Omega$.
  Por tanto, $G(v)\psi$ se factoriza a través de $\Omega_G$ por un morfismo, digamos, $v_G$.
  \begin{equation*}
    \begin{tikzcd}
      D\ar{r}{!_D}\ar[swap]{d}{m} & 1\ar{r}{\cong}\ar{d}{v_G} & G1\ar{d}{Gv} \\
      A\ar{r}{h^*}\ar[no head]{d}{} & \Omega_G\ar{r}{e} & G\Omega. \\
      {}\ar[no head]{rr}{G(h)s} & {} & {}\ar{u}{}
    \end{tikzcd}
  \end{equation*}
Ahora el cuadrado de la derecha es un producto fibrado al tener un isomorfismo arriba y un mono abajo,
mientras que el exterior es un producto fibrado al ser un rectángulo de~\eqref{eq:clasificadorm} (hay una única
flecha de $D$ a $G1\cong 1$), por tanto, el cuadrado de la izquierda es un producto fibrado. De esta manera,
para que $v_G\colon 1\to \Omega_G$ sea el clasificador de subobjetos en $\mathcal{E}_C$, solo falta mostrar que
$h^*$ es único. Supongamos que hay otro morfismo $f\colon (A,s)\to (\Omega_G,u)$ tal que el diagrama de la izquierda
es un producto fibrado. Para mostrar que $f=h^*$, basta mostrar que $ef=eh^*=G(h)s$, pero $G(h)s$ es un morfismo
hacia una coálgebra colibre, por lo que, pasando por la adjunción $U\dashv R$, es suficiente que
$\epsilon_{\Omega}ef=h$ en $\mathcal{E}$.
\begin{align*}
  \epsilon_{\Omega}ef &= \epsilon_{\Omega}G(\tau)\delta_{\Omega}ef \tag{Definición de e} \\
  &= \tau\epsilon_{G\Omega}\delta_{\Omega}ef \tag{Naturalidad de $\epsilon$} \\
  &= \tau ef. \tag{Estructura de comónada}
\end{align*}
Sin embargo tenemos el siguiente diagrama:
\begin{equation*}
  \begin{tikzcd}
    D\ar{r}{!_D}\ar[swap]{d}{m} & 1\ar{r}{\cong}\ar{d}{v_G} & G1\ar{r}{!_{G1}}\ar{d}{Gv} & 1\ar{d}{v} \\
    A\ar{r}{f} & \Omega_G\ar{r}{e} & G\Omega\ar{r}{\tau} & \Omega.
  \end{tikzcd}
\end{equation*}
Dado que $f$ clasifica a $m$, el cuadrado de la izquierda es un producto fibrado en $\mathcal{E}_C$, pero por
como construimos los límites finitos en un lema anterior, el funtor que olvida los preserva, entonces también es un
producto fibrado en $\mathcal{E}$. El cuadro de en medio tiene un iso arriba y un mono abajo, mientras que el cuadro
de la derecha es un producto fibrado al ser la definición de $\tau$. Lo anterior nos permite concluir que el exterior
es un producto fibrado y que $\tau ef$ clasifica a $m$ en $\mathcal{E}$, en otras palabras, $\tau ef = h$ que era
lo que necesitábamos para que $f=h^*$. Por lo tanto, $\mathcal{E}_G$ tiene clasificador de subobjetos y es un topos.
\end{proof}

\begin{proposition}
  Si $A$ es un espacio discreto, entonces $\Gamma A = A$.
\end{proposition}
\begin{proof}
  Por \ref{ax:gds7} existe un único subespacio discreto $\Gamma A$ de $A$ tal que todo punto de $A$ está en $\Gamma A$. 
  Dado que $A$ es discreto, $A$ es un subespacio de $A$, y cada punto de $A$ está en $A$, concluimos que $\Gamma A = A$.
\end{proof}

\begin{corollary}[label=cor:unidadidentidad]
  $\Gamma\Delta$ es la identidad en $M$ y $\Delta\Gamma$ es idempotente.
\end{corollary}

\begin{lemma}
  $\mathcal{C}=(\Delta\Gamma,c,=)$ es una comonada en Espacios.
\end{lemma}
\begin{proof}
  Para todo espacio $A$, tenemos que $\Delta\Gamma A = \Delta\Gamma\Delta\Gamma A$ por \ref{cor:unidadidentidad},
  lo cual se puede pensar como un isomorfismo natural $Id\colon \Delta\Gamma\to \Delta\Gamma\Delta\Gamma$, cuyas componentes son identidades. En esta prueba 
  $c$ es la counidad de la adjunción $\Delta\dashv\Gamma$, es decir, $c\colon \Delta\Gamma\to Id_{Espacios}$ que construimos en \ref{thm:funtordiscretos}.
  Procedemos a demostrar que se cumplen los diagramas de comonada para $\mathcal{C}$. Primero, la coasociatividad,
  \begin{equation*}
    \begin{tikzcd}
      \Delta\Gamma A\ar{r}{id_{\Delta\Gamma A}}\ar[swap]{d}{id_{\Delta\Gamma A}} & \Delta\Gamma\Delta\Gamma A\ar{d}{id_{\Delta\Gamma\Delta\Gamma A}} \\
      \Delta\Gamma\Delta\Gamma A\ar[swap]{r}{\Delta\Gamma id_{\Delta\Gamma A}} & \Delta\Gamma\Delta\Gamma\Delta\Gamma A.
    \end{tikzcd}
  \end{equation*}
  El diagrama conmuta gracias a que $\Delta\Gamma$ es functorial. Ahora, la counidad, queremos ver
  que el siguiente diagrama conmuta:
  \begin{equation*}
    \begin{tikzcd}
      {} & \Delta\Gamma A\ar{d}{id}\ar[swap]{ld}{id}\ar{rd}{id} & {} \\
      \Delta\Gamma A & \Delta\Gamma\Delta\Gamma A\ar{l}{c_{\Delta\Gamma A}}\ar[swap]{r}{\Delta\Gamma c_A} & \Delta\Gamma A.
    \end{tikzcd}
  \end{equation*}
  Dado que $u_{\Gamma A} = id_{\Gamma A}$, sucede que $id_{\Delta\Gamma A} = \Delta u_{\Gamma A}$, y
  ambos triángulos se sigue de las ecuaciones unidad-counidad, por lo que el diagrama conmuta.

\end{proof}

\begin{theorem}
  La categoría $M$ de espacios discretos es un topos.
\end{theorem}
\begin{proof}
  Basta demostrar que $M$ es equivalente al topos $Espacios_C$ de coálgebras, con
  $\mathcal{C}$ la comónada del lema anterior. Con tal propósito, proponemos al funtor
  $F\colon M\to Espacios_{\mathcal{C}}$ que envía a cada espacio discreto $D$ a la coálgebra 
  $(\Delta D, u_D)$, donde $u\colon Id\to\Delta\Gamma$ es la unidad de la adjunción, y en 
  morfismos $F$ es la identidad. Primero veamos que $(\Delta D, \Delta u_D)$ realmente es una coálgebra,
  para lo cual es necesario que los siguientes diagramas conmuten.
  \begin{equation*}
    \begin{tikzcd}
      \Delta D & \Delta\Gamma\Delta D\ar[swap]{l}{c_{\Delta D}} & {} & \Delta\Gamma\Delta\Gamma\Delta D & \Delta\Gamma\Delta D\ar[swap]{l}{\Delta\Gamma\Delta u_D} \\
       & \Delta D\ar{lu}{id}\ar[swap]{u}{\Delta u_D} & {} & \Delta\Gamma\Delta D\ar{u}{id} & \Delta D.\ar{l}{\Delta u_D}\ar[swap]{u}{\Delta u_D}
    \end{tikzcd}
  \end{equation*}
  El primer diagrama conmuta ya que es una ecuación unidad-counidad, mientras que
  el segundo diagrama se da gracias a que las componentes de la unidad son identidades.
  Además, dejar igual a los morfismos resulta en morfismos de coálgebras ya que si $f\colon D\to D'$
  es un morfismo de espacios discretos, entonces el siguiente diagrama conmuta:
  \begin{equation*}
    \begin{tikzcd}
      \Delta D\ar{r}{f}\ar[swap]{d}{\Delta u_D} & \Delta D'\ar{d}{\Delta u_{D'}} \\
      \Delta\Gamma\Delta D\ar[swap]{r}{\Delta\Gamma f} & \Delta\Gamma\Delta D'.
    \end{tikzcd}
  \end{equation*}
  Recordando que $\Delta$ es la inclusión, tenemos que el diagrama es un cuadrado de naturalidad de $u$.
  Por tanto, $F$ es funtorial ya que no hace nada en morfismos. Ahora, proponemos al funtor inverso $G\colon Espacios_{\mathcal{C}}\to M$ 
  que envía a cada coálgebra $(A,s)$ a $\Gamma A$ y a cada morfismo de coálgebras $f\colon (A,s)\to (B,t)$ a $\Gamma f$.
  Es fácil ver que $G$ es un funtor ya que $\Gamma$ lo es. Sea $D$ un espacio discreto y 
  $(A,s)$ una coálgebra, veamos qué sucede con las composiciones.
  \begin{align*}
    GF(D) &= G(\Delta D, \Delta u_D) = \Gamma\Delta D = D \\
    FG(A,s) &= F(\Gamma A) = (\Delta\Gamma A, \Delta u_{\Gamma A}) \cong (A,s).
  \end{align*}
  Para convencernos del último paso consideremos el siguiente diagrama:
  \begin{equation*}
    \begin{tikzcd}
      \Delta\Gamma A\ar[swap]{d}{\Delta u_{\Gamma A}}\ar{r}{c_A} & A\ar{d}{s} \\
      \Delta\Gamma\Delta\Gamma A\ar[swap]{r}{\Delta\Gamma c_A} & \Delta\Gamma A.
    \end{tikzcd}
  \end{equation*}
  Se trata de un diagrama conmutativo puesto que la parte inferior es la identidad, 
  una vez más, al ser una ecuación unidad-counidad. Luego, la condición de counidad para que
  $(A,s)$ sea coálgebra nos dice que $c_As = id_A$, por lo que $c_A$ es mono y tiene inverso por la derecha,
  así que $c_A$ es de hecho un isomorfismo con $s$ su inverso por ambos lados. Entonces la parte superior
  del diagrama también es la identidad, y por lo anterior, $c_A$ es un isomorfismo de coálgebras.
  Por lo tanto $F$ es una equivalencia de categorías y $M$ es un topos.
\end{proof}

\begin{theorem}
  Si $\Omega$ es el clasificador de subobjetos en Espacios, entonces $\Gamma\Omega$ con 
  $\Gamma v\colon \Gamma1\to\Gamma\Omega$ es el clasificador de subobjetos de M. $\Gamma$ preserva al 
  clasificador de subobjetos.
\end{theorem}
\begin{proof}
  Para cada $A\in M$ y $m\colon S\to A$ un subespacio de $A$, existe un único morfismo $\chi_m\colon A\to \Omega$ tal que el siguiente diagrama conmuta:
  \begin{equation*}
    \begin{tikzcd}
      \Delta S\ar{r}{!}\ar[swap,tail]{d}{\Delta m} & 1\ar{d}{v} \\
      \Delta A\ar[swap]{r}{\chi_m} & \Omega.
    \end{tikzcd}
  \end{equation*}
  Como $\Gamma$ es adjunto derecho, preserva límites, entonces al aplicar $\Gamma$ al diagrama anterior
  obtenemos un producto fibrado. Por \ref{prop:terminaldiscreto} $\Gamma 1 = 1$ y aplicamos \ref{cor:unidadidentidad}.
  \begin{equation*}
    \begin{tikzcd}
      S\ar{r}{!}\ar[swap,tail]{d}{m} & 1\ar{d}{\Gamma v} \\
      A\ar[swap]{r}{\Gamma \chi_m} & \Gamma\Omega.
    \end{tikzcd}
  \end{equation*}
  Si hubiera otro morfismo $f\colon A\to \Gamma\Omega$ tal que el diagrama sea un producto fibrado, entonces
  consideramos al siguiente diagrama:
  \begin{equation*}
    \begin{tikzcd}
      S\ar{r}{!}\ar[swap,tail]{d}{m} & \Delta\Gamma 1\ar{d}{\Delta\Gamma v}\ar{r}{c_1} & 1\ar{d}{v} \\
      A\ar[shift left=0.75ex]{r}{f}\ar[swap,shift right=0.75ex]{r}{\Gamma \chi_m} & \Delta\Gamma\Omega\ar[swap,tail]{r}{c_{\Omega}} & \Omega.
    \end{tikzcd}
  \end{equation*}
  El cuadrado de la derecha conmuta por la naturalidad de la counidad, pero es también un producto fibrado.
  \begin{equation*}
    \begin{tikzcd}
      X\ar[bend left=30]{rrd}{!}\ar[bend right=30]{rdd}[swap]{p}
          \ar[dashed]{rd}{!}\\[-2ex]
      &[-2ex] \Delta\Gamma 1\ar{r}{c_1}\ar[swap]{d}{\Delta\Gamma v} & 1\ar{d}{v} \\
      &\Delta\Gamma\Omega\ar[tail]{r}{c_{\Omega}} & \Omega.
    \end{tikzcd}
  \end{equation*}
  Bajo la hipótesis de que $v!=c_{\Omega}p$, como los diagramas que terminal en $1$ conmutan, 
  $c_{\Omega}p=vc_1!=c_{\Omega}\Delta\Gamma v!$, entonces $p=\Delta\Gamma v!$ y en efecto, es un producto fibrado.
  De esta manera, por el lema del producto fibrado, el exterior es un producto fibrado, lo cual implica que
  $c_{\Omega}f=\chi_m=c_{\Omega}\Gamma \chi_m$, pero $c_{\Omega}$ es mono, entonces $f=\Gamma \chi_m$. 
  Por tanto, $\Gamma\Omega$ con $\Gamma v$ es el clasificador de subobjetos de $M$.
\end{proof}

\begin{theorem}[label=thm:puntitos]
  Si $A$ tiene una cantidad finita de puntos, entonces $\Gamma A=\sqcup 1$ donde la cantidad de copias
  de $1$ es igual a la cantidad de puntos de $A$.
\end{theorem}
\begin{proof}
  Basta demostrarlo para $A$ con dos puntos, $a,b\colon 1\to A$ con $a\neq b$ y el resultado se sigue por inducción.
  Primero hay que ver que $\Gamma1 + \Gamma1$ es un subespacio de $A$. Como $\Gamma 1=1$ y tenemos elementos globales 
  globales $a$ y $b$, por la propiedad universal del coproducto, existe un morfismo $a+b\colon 1+1\to A$. Veamos que es un monomorfismo.
  \begin{equation*}
    \begin{tikzcd}
      (a+b)\times_A (a+b)\ar{r}\ar{d} & 1+1\ar{d}{a+b}\\
      1+1\ar{r}{a+b} & A.
    \end{tikzcd}
  \end{equation*}
  En un topos el coproducto es estable bajo productos fibrados, entonces $(a+b)\times_A (a+b)\cong
  (a\times_A a) + (a\times_A b) + (b\times_A a) + (b\times_A b)$. Además, $a\times_A b\cong 0$
  ya que $a\neq b$. Por lo tanto, $(a+b)\times_A (a+b)\cong (a\times_A a) + (b\times_A b)$, pero $a$ y $b$
  son monomorfismos, entonces $(a+b)\times_A (a+b)\cong 1+1$. Por lo anterior, $a+b$ es mono. Por otro lado,
  $\Gamma1+\Gamma1$ es un espacio discreto porque es el coproducto en $M$. Falta demostrar que los puntos
  de $A$ se factorizan a través de $\Gamma1+\Gamma1 = 1+1$.
  \begin{equation*}
    \begin{tikzcd}
      1\ar{r}{a}\ar[swap]{dr}{i} & \Gamma A & 1\ar{l}{b}\ar{dl}{i} \\
      {} & 1+1.\ar[dashed]{u}{a+b} 
    \end{tikzcd}
  \end{equation*}
  Los puntos de $A$ se encuentran en $\Gamma A$ y se factorizan a través de $a+b$. Por lo tanto,
  la unicidad de $\Gamma A$ implica que $\Gamma A = 1+1$.
\end{proof}

\begin{corollary}
  $M$ es un topos booleano.
\end{corollary}
\begin{proof}
  Sabemos por \ref{cor:bivaluado} que $\Omega$ tiene dos puntos, entonces por \ref{thm:puntitos} $\Gamma\Omega \cong 1+1$.
\end{proof}

\backmatter
\printbibliography


\end{document}